% ===================================================================
%  Remove This Dependency
%  Draft paper. Compile:  pdflatex remove-this-dependency.tex (x2)
% ===================================================================
\documentclass[10pt,conference]{IEEEtran}

\usepackage[utf8]{inputenc}
\usepackage[T1]{fontenc}
\usepackage{booktabs}
\usepackage{graphicx}
\usepackage{amsmath}
\usepackage{amssymb}
\usepackage{xcolor}
\usepackage{listings}
\usepackage{url}
\usepackage[hidelinks]{hyperref}
\usepackage{balance}

% ---- draft annotations: comment out \draftfalse to hide --------------
\newif\ifdraft
\drafttrue
\newcommand{\todo}[1]{\ifdraft\textcolor{red}{\textbf{[TODO: #1]}}\fi}
\newcommand{\pilot}[1]{\textcolor{blue}{#1}}

% ---- shorthands ------------------------------------------------------
\newcommand{\tool}{\textsc{Unpin}}
\newcommand{\bench}{\textsc{UnpinBench}}
\newcommand{\dibench}{DI-Bench}
\newcommand{\code}[1]{\texttt{\small #1}}

\lstset{
  basicstyle=\ttfamily\footnotesize,
  breaklines=true,
  columns=fullflexible,
  frame=single,
  framesep=3pt,
  xleftmargin=4pt,
  showstringspaces=false,
}

\begin{document}

\title{Remove This Dependency:\\
Why Passing Tests Cannot Tell You a Dependency Is Gone}

\author{
\IEEEauthorblockN{Anonymous Author(s)}
\IEEEauthorblockA{Anonymous Institution\\
\texttt{anonymous@example.edu}}
}

\maketitle

% =====================================================================
\begin{abstract}
Execution-based benchmarks judge a repository-level change by whether the
project's test suite still passes. We show this oracle is substantially
blind for dependency-related changes, and we quantify the blindness.
Mutating the ground-truth manifests of \dibench's Python subset one
dependency at a time, we find that \textbf{49.4\%} of declared
dependencies can be deleted without any test failing: 91 of 115 because
the package is still installed transitively (a \emph{phantom}), and 24
because no test reaches the code that uses it (a \emph{blind spot}).
Both are invisible to execution pass rate, the headline metric of the
benchmark. We then show the blindness is predictable without running
anything: a two-layer static signal over the dependency graph and the
import graph identifies CI-invisible deletions with precision 0.81 and
recall 0.83.

Building on this, we introduce \bench, an executable benchmark for a task
that existing work does not cover: removing a dependency that the code
\emph{actually uses}, by rewriting the code, with behaviour preserved.
Unlike dependency-removal datasets mined from commit histories, every
\bench{} instance is verified so that deleting the dependency provably
breaks CI, which makes the test suite a meaningful oracle for that
instance. Because tests alone still accept several ways of faking a
removal, we pair the benchmark with a six-gate oracle and validate it
against 176 automatically constructed cheats. Tests alone accept 35.8\%
of cheats; install-level gates close vendoring and hiding completely,
while 8 stubbed removals pass every install-level check and are caught
only by a structural signal.

A pilot of two open-weight coding agents over 28 runs finds that
execution pass rate reports 5 successes where only 3 are genuine
removals. We release the benchmark, the harness, which runs on free
continuous-integration infrastructure, and all artifacts.
\end{abstract}

\begin{IEEEkeywords}
software dependencies, benchmarks, coding agents, test oracles,
mutation analysis, program analysis
\end{IEEEkeywords}

% =====================================================================
\section{Introduction}

A program's dependencies are a property of its code, not of its manifest
file. The manifest is a claim about that property, and claims drift.
Drosos et al.~\cite{drosos2024bloat} found that over half of declared
dependencies in the PyPI ecosystem are bloated, and that the single
largest cause---34\% of cases---is refactoring that removed the last use
of a library without removing its declaration. Code and build
configuration co-evolve, and they co-evolve imperfectly: McIntosh et
al.~\cite{mcintosh2011build} reported that up to a quarter of source
changes and nearly half of test changes require a matching build change.

Benchmarks that evaluate automated systems on dependency tasks resolve
this by execution. \dibench~\cite{dibench2025} masks a repository's
dependency declarations, asks a model to restore them, and scores the
result by whether the project's continuous-integration (CI) suite passes.
This is an appealing oracle: it is objective, it is automatic, and it
measures something users care about.

This paper argues that it is also substantially blind, and quantifies how
blind. The argument is simple. A test suite can only notice a missing
dependency if some test reaches the code that imports it. A test suite
can only notice an \emph{extra} dependency if that dependency fails to
install. Neither condition holds often enough for the oracle to be
trusted on its own.

\paragraph*{Contributions}
\begin{enumerate}
  \item \textbf{A measurement of oracle blindness
  (\S\ref{sec:blindness}).} We mutate the ground-truth manifests of
  \dibench's Python subset, deleting one dependency at a time and
  re-running the real CI suite in a container: 531 mutants in total.
  On instances whose ground-truth patch still passes, 49.4\% of
  deletions go unnoticed. We separate the two mechanisms---phantom
  installation and test blind spots---because they call for different
  remedies.

  \item \textbf{A static predictor of that blindness
  (\S\ref{sec:predictor}).} Combining the package-level dependency graph
  with import-level reachability predicts which deletions CI will miss,
  at precision 0.81 and recall 0.83, with no container and no test run.

  \item \textbf{\bench{} (\S\ref{sec:benchmark}).} An executable
  benchmark for in-situ dependency removal, where each instance is
  verified so that deleting the target provably fails CI. We describe
  how instances are mined and verified, and release a pool of 162
  candidate instances over 47 repositories.

  \item \textbf{A gate stack validated against constructed cheats
  (\S\ref{sec:gates}).} Six gates, of which tests are one. We generate
  176 cheats automatically---vendoring, stubbing, and hiding---and show
  that tests alone accept 35.8\% of them, that install-level gates close
  two cheat families entirely, and that stubbing survives every
  install-level check.

  \item \textbf{An agent pilot (\S\ref{sec:agents})} over two
  open-weight coding agents and 28 runs, in which execution pass rate
  reports 5 successes against 3 genuine removals.
\end{enumerate}

\paragraph*{Scope} All experiments are on Python. \S\ref{sec:threats}
discusses what this costs us. The harness, the mutants, the cheat
generator, and the agent scaffold are released; the harness runs on free
GitHub Actions runners, so every result here is reproducible without
dedicated hardware.

% =====================================================================
\section{Background and Related Work}
\label{sec:related}

\subsection{Dependency inference and its oracle}

\dibench~\cite{dibench2025} is the closest benchmark to our setting: 581
repositories across four languages, dependency declarations masked, and
both textual metrics (precision and recall over package names) and an
execution metric. Its authors note that ``test coverage for each
repository may not be exhaustive'' and leave the consequence unexplored.
We quantify that consequence.

DependEval~\cite{dependeval2025} evaluates repository dependency
\emph{understanding} through static tasks and does not execute.
Environment-setup benchmarks---EnvBench~\cite{envbench2025},
Repo2Run~\cite{repo2run2025}, ExecutionAgent~\cite{executionagent2025}---
ask an agent to make a repository runnable at all, a different task with
the same oracle and, we expect, the same blindness.

\subsection{Tests as an oracle}

That test suites miss dependency-induced faults is known. Hejderup and
Gousios~\cite{hejderup2022trust} seeded faults into library call sites in
Java projects and found test suites detected a minority. More recently,
TimeMachine-bench~\cite{timemachine2026} reports agents producing
``spurious solutions that exploit low test coverage'' on dependency
upgrades, and SWE~Atlas~\cite{sweatlas2026} measures a 15--40 point gap
between passing behavioural checks and meeting maintainability rubrics on
refactoring tasks. Our contribution is to mutate the \emph{manifest}
rather than the code, which isolates the dependency oracle specifically,
and to separate phantom installation from genuine test blindness.

\subsection{Dependency bloat and debloating}

A substantial literature detects and removes \emph{unused} dependencies:
DepClean~\cite{soto2021longitudinal} for Maven,
PyTrim~\cite{pytrim2025} for Python, and the fine-grained PyPI analysis
of Drosos et al.~\cite{drosos2024bloat}. These tools solve the easy case,
where no code calls the library. \bench{} targets the hard case, where
the library is in use and removal requires writing a replacement.

\subsection{Removal and migration datasets}

He et al.~\cite{he2021migration} mined 19,652 Java projects and found
8,065 with at least one library removal; a 2024 dataset of
dependency-removal commits and their pull requests is publicly
available~\cite{jaisri2024dataset}. Both are \emph{mined metadata}: they
record that removals happened, but cannot score a candidate removal,
because they carry no executable environment. Library \emph{migration}
benchmarks---PyMigBench~\cite{pymigbench2023},
MigrationBench~\cite{migrationbench2025}---cover replacing library $A$
with library $B$, not eliminating a library in favour of one's own code.
LibReuseBench~\cite{libreuse2026} and Commit0~\cite{commit02024} ask
models to write code that avoids or reimplements libraries, but on
greenfield problems with no existing code whose behaviour must be
preserved. To our knowledge no executable benchmark exists for in-situ
removal of a used dependency.

\subsection{Agent scaffolds and reward hacking}

Our agent scaffold follows the agent-computer interface of
SWE-agent~\cite{sweagent2024}: a line-numbered viewer with a 100-line
window, search results capped at 50, a range-based edit command with an
integrated syntax linter, and collapsing of observations older than the
last five. Evaluation uses greedy decoding at temperature 0 with a single
attempt, the standard \textit{pass@1} protocol.

That agents exploit weak verifiers is well documented.
SpecBench~\cite{specbench2026} separates visible from held-out tests and
finds most of the gap arises from compositional failure rather than
deliberate cheating; a hacker--fixer loop~\cite{hackerfixer2026} found
16\% of tasks across five agent benchmarks hackable from the task
description alone. We do not claim to measure reward hacking in general.
We claim something narrower and testable: on one concrete task, we
enumerate the ways a change can pass the test oracle without doing the
work, and we measure which gates catch which.

% =====================================================================
\section{How Blind Is the Test Oracle?}
\label{sec:blindness}

\subsection{Setup}

We use the Python subset of \dibench{} (98 repositories). Each instance
carries a masked manifest, a ground-truth (``gold'') patch restoring the
declarations, and an \code{act} command that replays the repository's own
GitHub Actions workflow inside a container. We re-implemented the
execution harness to run on free GitHub Actions runners under the
\code{sysbox} runtime\footnote{The published harness requires
\code{sysbox}, which is unavailable on standard CI runners; we install it
in the job. Plain privileged Docker-in-Docker fails because the inner
daemon cannot stack overlay filesystems.}.

\paragraph*{Benchmark decay}
Of 98 gold patches, only \textbf{51 of 96} still pass CI at the time of
writing, roughly 21 months after release.\footnote{Two instances are
excluded: their \code{act} job never started, because an action they use
requires a newer Node runtime than \code{act} supports. \code{act} prints
\emph{Job succeeded} in this situation, so an evaluator that trusts its
output alone records a false pass.} The failures are environmental drift:
an installer that dropped support for old Python versions, a discontinued
distribution download, and unpinned libraries whose newer releases break
the projects' own assertions. This is a measurement problem for anyone
using the benchmark today, and it is the same drift the paper is about,
appearing inside the benchmark itself. \emph{All execution numbers below
are therefore conditioned on gold-passing instances.}

\subsection{Dependency mutation analysis}

For every declared dependency $d$ of every instance, we construct a
mutant: the gold manifest with $d$'s declaration removed and nothing else
changed. We verify syntactically that the resulting dependency set equals
the gold set minus $d$. Each mutant is then scored by the real harness.
This yields 531 mutants, of which 233 belong to gold-passing instances.

We call the fraction of deletions that still pass CI the
\emph{dependency mutation score}. It bounds what execution pass rate can
ever detect: a dependency whose deletion goes unnoticed cannot be scored
correctly by an execution oracle, whether the system under test omits it
or invents it.

\begin{table}[t]
\caption{Dependency mutation analysis, \dibench{} Python. Mutants of
gold-passing instances.}
\label{tab:mutation}
\centering
\begin{tabular}{lr}
\toprule
\textbf{Measure} & \textbf{Value} \\
\midrule
Deletion mutants scored                       & 233 \\
Deletions that still pass CI                  & 115 (49.4\%) \\
\quad of which \emph{phantom} (installed anyway) & 91 \\
\quad of which \emph{blind spot} (absent, tests pass) & 24 (10.3\%) \\
Instances with $\geq 1$ invisible dependency  & 38 of 51 \\
Instances with $\geq 1$ blind-spot dependency & 15 of 51 \\
\midrule
Sanity: visible deletions where the package & \\
\quad was indeed absent from the environment  & 114 of 118 \\
\bottomrule
\end{tabular}
\end{table}

\subsection{Two mechanisms, two remedies}

Table~\ref{tab:mutation} separates the result by mechanism, determined by
parsing the CI installation log for each mutant.

\textbf{Phantom installation (91 cases).} The declaration is removed, but
the package is installed regardless, pulled in transitively by another
declared dependency or by a test-runner configuration. The code's imports
resolve and the suite passes. These declarations are redundant: this is
the ``transitive dependency listed unnecessarily'' category that Drosos
et al.\ measure at 30\% of bloat causes~\cite{drosos2024bloat}, here
sitting inside a benchmark's ground truth. \dibench's curation verified
that removing \emph{all} dependencies fails; it never verified that
removing \emph{one} does.

\textbf{Test blind spot (24 cases, 10.3\% of mutants).} The package is
genuinely absent from the environment, the import would fail, and the
suite still passes, because no test reaches the importing code. These are
real coverage gaps, present in 15 of 51 instances.

The remedies differ. Phantoms are a ground-truth quality problem, fixable
by curation. Blind spots are intrinsic to the projects' test suites and
can only be addressed by adding an oracle that does not depend on tests.

% =====================================================================
\section{Predicting Blindness Without Execution}
\label{sec:predictor}

If an execution oracle is blind in a predictable way, the prediction is
itself useful: for curating instances, for flagging unreliable scores,
and for building the non-test gates of \S\ref{sec:gates}. We test whether
two static signals suffice.

\textbf{Dependency layer.} Package $d$ is predicted phantom if it lies in
the transitive requirement closure of the \emph{other} declared
dependencies, computed from PyPI metadata (\code{requires\_dist},
ignoring extras-conditioned requirements, depth~6).

\textbf{Code layer.} Package $d$ is predicted unreachable if no test file
reaches, directly or transitively through project modules, a file that
imports it.

\begin{table}[t]
\caption{Static prediction of CI-invisible deletions, evaluated against
the 233 executed mutants (115 invisible).}
\label{tab:predictor}
\centering
\begin{tabular}{lrrr}
\toprule
\textbf{Predictor} & \textbf{Pred.} & \textbf{Prec.} & \textbf{Rec.} \\
\midrule
Dependency layer only  & 73  & 0.97 & 0.62 \\
Code layer only        & 70  & 0.69 & 0.42 \\
Both (disjunction)     & 118 & 0.81 & 0.83 \\
\bottomrule
\end{tabular}
\end{table}

Table~\ref{tab:predictor} reports the result. The dependency layer is
almost never wrong, which matches its mechanism: if the package is in the
closure it \emph{will} be installed. The code layer is noisier, because
static import reachability under-approximates dynamic imports. Together
they recover five in six invisible deletions at four in five precision,
from the manifest, public package metadata, and import statements alone.

\paragraph*{Symbol-level reachability} We also evaluated a symbol-level
variant using PyCG~\cite{pycg2021} call graphs. PyCG completed on 18 of
51 repositories within a four-minute budget and timed out on the rest; on
those it covered, symbol-level reachability performed \emph{worse} than
the file-level signal (precision 0.47, recall 0.64), because missed
dynamic calls cause it to report reachable dependencies as unreachable.
We therefore use the file-level signal and leave a scalable symbol-level
layer to future work.

% =====================================================================
\section{\bench: Removing a Dependency That Is Used}
\label{sec:benchmark}

\subsection{Task definition}

An instance is a triple $(R, d, T)$: a repository $R$ at a fixed commit,
a target package $d$ that $R$'s code imports, and $R$'s own test suite
$T$. The task is to produce a change to $R$ such that $R$ no longer
requires $d$, with behaviour preserved. The agent must rewrite the code
that uses $d$, remove $d$'s declaration, and leave $T$ passing.

This is deliberately the \emph{hard} case. Removing an unused declaration
is solved by static tools~\cite{pytrim2025,soto2021longitudinal}.
Migrating from $A$ to $B$ is covered by existing
benchmarks~\cite{pymigbench2023,migrationbench2025}. Removing a library
that is in use, by writing the replacement, is neither, and it is what
developers do when a dependency is abandoned, vulnerable, or
disproportionate to what it provides.

\subsection{Instance verification}

The defining property of a \bench{} instance is that \emph{deleting $d$'s
declaration provably fails CI}. This is exactly the mutation result of
\S\ref{sec:blindness} used constructively: we keep only those
$(R,d)$ pairs whose deletion mutant \emph{failed}. Without this filter,
an instance would be unscoreable---an agent could delete the declaration,
change nothing, and pass.

This verification step is what distinguishes \bench{} from mined removal
datasets~\cite{jaisri2024dataset,he2021migration}. Those record that a
removal occurred; they cannot tell whether a candidate removal is
correct, because they carry no executable environment.

\subsection{Instance pool and difficulty}

From the gold-passing \dibench{} Python instances we derive 162 candidate
pairs over 47 repositories, in three tiers:

\begin{itemize}
  \item \textbf{Strong} (62 pairs, 38 repositories): the package is
  imported by one or two source files, is reachable from tests, and its
  deletion fails CI. These are scoreable and tractable.
  \item \textbf{Medium} (19 pairs): three to five importing files.
  \item \textbf{Weak} (81 pairs): deletion does \emph{not} fail CI. These
  are excluded from scoring but retained, since they are precisely the
  instances where the test oracle is blind.
\end{itemize}

We label each instance by \emph{used API surface}: the number of distinct
symbols of $d$ that $R$ calls, plus call sites. In the strong tier this
ranges from 1 to 22 (median 1). We hypothesise, and \S\ref{sec:agents}
begins to test, that success falls and cheating rises with surface size.

\todo{full run: extend the pool by mining removal commits (repositories
cloned with full history) to obtain human reference solutions, and extend
to \dibench{}'s \emph{large} subset (50 further Python repositories).}

% =====================================================================
\section{A Gate Stack, and What Tests Miss}
\label{sec:gates}

\subsection{Six gates}

Tests are one oracle among several. We define six gates; a change is
accepted only if it passes all of them.

\begin{description}
  \item[G1 Declaration gone.] $d$ is absent from the manifest and from
  the resolved environment.
  \item[G2 Tests pass.] $R$'s own CI suite passes in a container.
  \item[G3 No vendored copy.] No file added by the change is
  token-similar to $d$'s own distributed source.
  \item[G4 Behaviour preserved beyond the tests.] The replacement is not
  a stub: added functions have non-trivial bodies, and no added package
  shadows $d$'s import name.
  \item[G5 No phantom use.] No import is now served by a transitively
  installed package.
  \item[G6 Quality not worse.] The replacement does not degrade
  structural quality.
\end{description}

G1 and G2 are what an execution benchmark already does. G3--G6 are the
contribution. \todo{full run: G6 is specified but not yet instantiated;
we plan to use a learned pairwise refactoring-quality metric.}

\subsection{Validating the gates against constructed cheats}

A gate stack is only as good as the attacks it is tested against. For
each of the 62 strong-tier instances we generate three cheats
automatically:

\begin{description}
  \item[Hide] move $d$ from the main dependency list to an optional
  extra or development group, which CI installs anyway;
  \item[Vendor] copy $d$'s own distributed source into the repository and
  delete the declaration;
  \item[Stub] create a local package shadowing $d$'s import name, whose
  imported symbols are no-ops, and delete the declaration.
\end{description}

Each variant is a patch against the masked repository, scored by the same
harness as any other candidate. Of 186 possible variants, 176 were
generated; 10 vendor variants were skipped because the package ships
compiled extensions and cannot be vendored as source.

\begin{table}[t]
\caption{Constructed cheats through the real CI harness. ``Caught''
counts variants that pass G2 but fail another gate.}
\label{tab:cheats}
\centering
\begin{tabular}{lrrrl}
\toprule
\textbf{Cheat} & \textbf{$n$} & \textbf{Pass G2} & \textbf{Caught} & \textbf{By} \\
\midrule
Hide   & 62 & 35 & 35 of 35 & G1, G5 \\
Vendor & 52 & 18 & 18 of 18 & G3 \\
Stub   & 62 &  8 &  0 of 8  & --- \\
\midrule
Total  & 176 & 61 (34.7\%) & 53 of 61 & \\
\bottomrule
\end{tabular}
\end{table}

Table~\ref{tab:cheats} gives the result, and it is the central
justification for the gate stack. Tests alone accept 61 of 176 cheats.
Install-level gates close hiding and vendoring \emph{completely}: every
cheat of those families that passes the tests fails G1, G5, or G3.

Stubbing is different. Eight stub variants pass CI, and no install-level
gate sees them, because nothing is installed and nothing is copied. Their
replacement code is between 4 and 27 lines of functions returning
nothing. They are detectable only structurally: our trivial-body and
shadow-package signals flag 61 of 62 stub variants. This is why G4 is
necessary rather than decorative, and it is the sharpest statement we can
make about what an execution oracle cannot do.

% =====================================================================
\section{Agent Pilot}
\label{sec:agents}

\paragraph*{Status} This section reports a \pilot{pilot} sized to a
\$3 budget. It establishes feasibility and direction; it does not support
rate estimates. \todo{full run: frontier models, all 62 strong instances,
three seeds, and the strict-prompt condition as an ablation.}

\subsection{Setup}

We evaluate two open-weight coding models through an OpenRouter endpoint:
\code{deepseek-v3.2} and \code{qwen3-coder}, chosen as the strongest
models within budget. The scaffold follows SWE-agent's interface
(\S\ref{sec:related}) with a 32-step budget. The task prompt is
deliberately \emph{neutral}: it states the goal and an expected procedure
but names no prohibition, so that any cheat observed is unprompted.
16 strong-tier instances were selected, nine with the largest used API
surface and seven with the smallest; both models stopped at a per-model
budget guard, yielding 15 and 13 runs respectively.

\subsection{Results}

\begin{table}[t]
\caption{Agent pilot, 28 runs. ``Genuine'' removals pass every gate.}
\label{tab:agents}
\centering
\begin{tabular}{lrrrr}
\toprule
\textbf{Model} & \textbf{Runs} & \textbf{Pass G2} & \textbf{Genuine} & \textbf{Cheats} \\
\midrule
\code{deepseek-v3.2} & 15 & 5 & 3 & 0 \\
\code{qwen3-coder}   & 13 & 0 & 0 & 0 \\
\midrule
Total & 28 & 5 & 3 & 0 \\
\bottomrule
\end{tabular}
\end{table}

Table~\ref{tab:agents} reports three findings.

\textbf{The test oracle overstates success by 67\%.} CI reports five
passes; three are genuine removals. Of the other two, one left the
declaration in place while CI passed anyway, and one changed nothing at
all, the package remaining installed transitively. Both are caught by
G1 and G5. On a task where the headline number is 5, the correct number
is 3.

\textbf{No spontaneous cheating was observed.} Across 28 runs there is no
vendoring, no stubbing, and no weakened test suite. The constructed
cheats of \S\ref{sec:gates} are therefore not an observation about these
models; they are a statement about what the oracle would accept. The
closest case is an agent that replaced \code{wcwidth} with a
self-described ``simplified replacement'' for a Unicode width function,
which passes CI and is plausibly a silent behaviour change.

\textbf{The task is hard.} Three clean successes in 28 runs. Failures are
mechanical as often as semantic: four agents left the manifest
syntactically invalid, others produced import errors or undefined names.

\paragraph*{Interpretation} The honest reading is that these models fail
this task far more often than they fake it. This does not weaken the case
for the gate stack---\S\ref{sec:gates} shows the gates are necessary
against attacks the oracle admits---but it does bound what the pilot
shows. Whether more capable models, which succeed more often, also cheat
more often is the question the full study must answer.

% =====================================================================
\section{Threats to Validity}
\label{sec:threats}

\paragraph*{Pilot scale} The agent results rest on 28 runs over 16
instances and two models, sized to a \$3 budget. Confidence intervals on
proportions at this scale span roughly $\pm$18 points; we report counts,
not rates. The mutation and cheat results, which carry the paper's main
claims, rest on 531 and 176 executed variants respectively.

\paragraph*{Step budget confound} \code{deepseek-v3.2} exhausted its step
budget on 13 of 15 runs, so its patches are frequently incomplete work
rather than considered answers. Its lower success rate partly reflects
our budget, not the model.

\paragraph*{Benchmark decay} Only 51 of 96 gold patches still pass. We
condition on gold-passing instances throughout, but the surviving set may
be systematically simpler, for instance better pinned.

\paragraph*{Harness fidelity} We found and report two defects in the
upstream evaluation path: the fake-package check crashes on VCS-URL
dependencies, and \code{act} prints \emph{Job succeeded} for a job that
never started, which an evaluator that ignores the exit code records as a
pass. We fixed both; results predating the fix are excluded. Measurement
infrastructure requires auditing as much as the systems it measures.

\paragraph*{Static predictor approximations} The dependency layer ignores
extras-conditioned requirements and installs driven by CI configuration
rather than the manifest, which is the main source of its missed
phantoms. The code layer under-approximates dynamic imports.

\paragraph*{Language scope} All results are Python. The mechanisms---
transitive installation and uneven test coverage---are not
Python-specific, but their magnitudes certainly are: ecosystems with
stricter dependency resolution should exhibit fewer phantoms.

\paragraph*{Cheat realism} Our cheats are constructed, not observed. They
establish what the oracle \emph{admits}, not what agents \emph{do}.

% =====================================================================
\section{Conclusion}

Execution-based benchmarks are trusted because running code seems
objective. We show that for dependency tasks this trust is misplaced in a
measurable way: in \dibench's Python subset, half of all declared
dependencies can be deleted with no test noticing. We give a static
predictor of that blindness, an executable benchmark for the removal
task whose instances are verified so the oracle is meaningful, and a gate
stack validated against constructed attacks, where tests alone accept a
third of cheats and stubbing defeats every install-level check. A pilot
of two open-weight agents finds execution pass rate reporting five
successes against three genuine removals.

The broader point is that an oracle should be measured, not assumed.
Mutation analysis of the ground truth is cheap, mechanical, and applicable
to any execution-based benchmark. We suspect many would not survive it.

% =====================================================================
\section*{Availability}

The harness, mutants, cheat generator, agent scaffold, gate
implementations, and all result artifacts are available at
\url{https://github.com/ANONYMIZED/di-bench-eval}. The harness runs on
free GitHub Actions runners.

% =====================================================================
\balance
\begin{thebibliography}{99}
\footnotesize

\bibitem{dibench2025}
Microsoft, ``DI-Bench: Benchmarking Large Language Models on Dependency
Inference with Testable Repositories at Scale,'' \emph{arXiv:2501.13699},
2025.

\bibitem{drosos2024bloat}
G.-P. Drosos, T. Sotiropoulos, D. Spinellis, and D. Mitropoulos,
``Bloat beneath Python's Scales: A Fine-Grained Inter-Project Dependency
Analysis,'' in \emph{Proc. FSE}, 2024.

\bibitem{mcintosh2011build}
S. McIntosh, B. Adams, T. H. D. Nguyen, Y. Kamei, and A. E. Hassan,
``An Empirical Study of Build Maintenance Effort,'' in \emph{Proc. ICSE},
2011.

\bibitem{hejderup2022trust}
J. Hejderup and G. Gousios, ``Can We Trust Tests to Automate Dependency
Updates? A Case Study of Java,'' \emph{Empirical Software Engineering},
2022.

\bibitem{timemachine2026}
``TimeMachine-bench: A Benchmark for Evaluating Model Capabilities in
Repository-Level Migration Tasks,'' in \emph{Proc. EACL}, 2026.

\bibitem{sweatlas2026}
``SWE Atlas: Benchmarking Coding Agents Beyond Issue Resolution,'' 2026.

\bibitem{soto2021longitudinal}
C. Soto-Valero, T. Durieux, and B. Baudry, ``A Longitudinal Analysis of
Bloated Java Dependencies,'' in \emph{Proc. ESEC/FSE}, 2021.

\bibitem{pytrim2025}
``PyTrim: A Practical Tool for Reducing Python Dependency Bloat,'' 2025.

\bibitem{he2021migration}
H. He, R. He, H. Gu, and M. Zhou, ``A Large-Scale Empirical Study on Java
Library Migrations,'' in \emph{Proc. ESEC/FSE}, 2021.

\bibitem{jaisri2024dataset}
P. Jaisri, ``The commit history of the dependent libraries and the
associated pull requests related to dependency removal,'' Zenodo,
\texttt{10.5281/zenodo.12730734}, 2024.

\bibitem{pymigbench2023}
M. Islam, A. A. Jahan, et al., ``PyMigBench: A Benchmark for Python
Library Migration,'' in \emph{Proc. MSR}, 2023.

\bibitem{migrationbench2025}
Amazon Science, ``MigrationBench: Repository-Level Code Migration
Benchmark from Java 8,'' \emph{arXiv:2505.09569}, 2025.

\bibitem{libreuse2026}
``LibReuseBench: A Test-Backed Benchmark for Library Reuse,'' 2026.

\bibitem{commit02024}
W. Zhao, N. Jiang, et al., ``Commit0: Library Generation from Scratch,''
\emph{arXiv:2412.01769}, 2024.

\bibitem{dependeval2025}
``DependEval: Benchmarking LLMs for Repository Dependency
Understanding,'' in \emph{Findings of ACL}, 2025.

\bibitem{envbench2025}
A. Eliseeva et al., ``EnvBench: A Benchmark for Automated Environment
Setup,'' 2025.

\bibitem{repo2run2025}
``Repo2Run: Automated Building Executable Environment for Code
Repositories,'' 2025.

\bibitem{executionagent2025}
I. Bouzenia and M. Pradel, ``ExecutionAgent: Automatic Test Execution
Setup,'' 2025.

\bibitem{sweagent2024}
J. Yang, C. E. Jimenez, A. Wettig, et al., ``SWE-agent: Agent-Computer
Interfaces Enable Automated Software Engineering,'' in \emph{Proc.
NeurIPS}, 2024.

\bibitem{specbench2026}
``SpecBench: Measuring Reward Hacking in Long-Horizon Coding Agents,''
2026.

\bibitem{hackerfixer2026}
Z. Zhong, I. Segal, et al., ``Hardening Agent Benchmarks with Adversarial
Hacker-Fixer Loops,'' 2026.

\bibitem{pycg2021}
V. Salis, T. Sotiropoulos, P. Louridas, D. Spinellis, and D. Mitropoulos,
``PyCG: Practical Call Graph Generation in Python,'' in \emph{Proc.
ICSE}, 2021.

\end{thebibliography}

\end{document}
